% ===================== Chapter 22 =====================
\chapter{Bilingual Programming in Salam}
\label{ch:bilingual}

We have mentioned throughout this book that Salam can be written in Persian as
well as English. Now we give that feature the chapter it deserves. Salam's
bilingualism is not a translation layer bolted on top --- it is built into the
language: the very same compiler reads English and Persian keywords, and produces
identical programs from either. This chapter shows how, and why it matters.

\section{One language, two vocabularies}

The key idea is simple. A programming language is made of \emph{keywords} (the
fixed words like \code{func}, \code{if}, \code{while}) and \emph{your} names (the
variables and functions you invent). Salam keeps the grammar --- the structure,
the rules, the meaning --- exactly the same in both languages, and changes only the
spelling of the keywords. \code{func} becomes \code{\fa{تابع}}; \code{if} becomes
\code{\fa{اگر}}; everything else about the program is unchanged.

Here is the same program twice. First in English:

\begin{salam}
func add(a i32, b i32) i32:
    ret a + b
end

func main:
    mut total i32 = 0
    for mut i = 1; i <= 5; i = i + 1:
        total = total + i
    end
    println "sum 1..5 =", total
end
\end{salam}

And now, exactly the same program, in Persian:

\begin{salamfa}
تابع جمع(الف صحیح۳۲, ب صحیح۳۲) صحیح۳۲:
    بازگشت الف + ب
پایان

تابع اصلی:
    متغیر مجموع صحیح۳۲ = 0
    برای متغیر i = 1; i <= 5; i = i + 1:
        مجموع = مجموع + i
    پایان
    چاپ "مجموع ۱ تا ۵ =", مجموع
پایان
\end{salamfa}

\begin{progout}
مجموع ۱ تا ۵ = 15
\end{progout}

Look closely and you will see the structures are identical, line for line. Only
the keywords and the type names changed. A Persian-speaking student reads the
second version as naturally as you read the first --- and they are learning exactly
the same concepts, because they \emph{are} the same concepts.

\section{The keyword equivalents}

Every keyword you have learned has a Persian twin. Here are the most common; the
complete list is in Appendix~A.

\begin{center}
\begin{tabular}{@{}llll@{}}
\toprule
\textbf{English} & \textbf{Persian} & \textbf{English} & \textbf{Persian} \\
\midrule
\code{func}   & \code{\fa{تابع}}    & \code{ret}     & \code{\fa{بازگشت}} \\
\code{if}     & \code{\fa{اگر}}     & \code{else}    & \code{\fa{وگرنه}} \\
\code{while}  & \code{\fa{تاوقتی}}  & \code{for}     & \code{\fa{برای}} \\
\code{repeat} & \code{\fa{تکرار}}   & \code{to}      & \code{\fa{تا}} \\
\code{mut}    & \code{\fa{متغیر}}   & \code{const}   & \code{\fa{ثابت}} \\
\code{struct} & \code{\fa{ساختار}}  & \code{enum}    & \code{\fa{شمارش}} \\
\code{interface} & \code{\fa{واسط}} & \code{end}     & \code{\fa{پایان}} \\
\code{import} & \code{\fa{واردکردن}} & \code{pub}    & \code{\fa{عمومی}} \\
\code{true}   & \code{\fa{درست}}    & \code{false}   & \code{\fa{نادرست}} \\
\code{break}  & \code{\fa{بشکن}}    & \code{continue} & \code{\fa{ادامه}} \\
\code{println} & \code{\fa{چاپ}}    & \code{print}   & \code{\fa{بنویس}} \\
\code{main}   & \code{\fa{اصلی}}    & \code{this}    & \code{\fa{این}} \\
\bottomrule
\end{tabular}
\end{center}

The type names have Persian forms too, which you first saw in Chapter~2:

\begin{center}
\begin{tabular}{@{}llll@{}}
\toprule
\textbf{English} & \textbf{Persian} & \textbf{English} & \textbf{Persian} \\
\midrule
\code{i32} & \code{\fa{صحیح۳۲}} & \code{i64} & \code{\fa{صحیح۶۴}} \\
\code{f64} & \code{\fa{اعشاری۶۴}} & \code{bool} & \code{\fa{منطقی}} \\
\code{str} & \code{\fa{رشته}}   & \code{char} & \code{\fa{کاراکتر}} \\
\code{auto} & \code{\fa{خودکار}} & & \\
\bottomrule
\end{tabular}
\end{center}

\begin{note}
The entry point itself is part of the vocabulary: it is \code{main} in English and
\code{\fa{اصلی}} in Persian. When you run a Persian program, the compiler looks for
\code{\fa{اصلی}} as the starting point, just as it looks for \code{main} in an
English one.
\end{note}

\section{Compiling a Persian program}

A Salam source file is written in \emph{one} language, and you tell the compiler
which with the \code{--lang} flag. English is the default; for a Persian file, pass
\code{--lang=fa}:

\begin{shell}
salam run hello.salam --lang=fa
salam build hello.salam --lang=fa --output=hello
\end{shell}

Everything else about the workflow --- \code{run}, \code{build}, \code{exec},
\code{fmt} --- is identical. The flag simply tells the compiler which keyword
vocabulary to expect.

\begin{warn}
A single file uses a single language: you cannot mix English and Persian
\emph{keywords} in the same file. A file is either an English program or a Persian
one, compiled with the matching \code{--lang}. This is deliberate --- it keeps each
file coherent and unambiguous. What you \emph{can} freely mix is covered next.
\end{warn}

\section{What you can mix}

The keyword rule is about keywords only. The \emph{names you choose} --- variables,
functions, struct fields --- may be in whatever language you like, in either kind of
file, because Salam identifiers may contain Unicode letters. So a Persian program
can use English identifiers (as the loop variable \code{i} did above), and an
English program can use Persian identifiers:

\begin{salam}
func main:
    naam str = "Sara"        // a Latin-letter Persian word as a name
    println "Hello,", naam
end
\end{salam}

\begin{progout}
Hello, Sara
\end{progout}

Likewise, \emph{strings} and \emph{comments} are free text and may be in any
language regardless of the file's keyword language --- which is why the Persian
program above could print the Persian message \code{"\fa{مجموع ۱ تا ۵ =}"} and an
English program can hold Persian comments. Keywords pick a side; your content does
not have to.

\begin{note}
Notice the digits. Inside the type name \code{\fa{صحیح۳۲}} the digits are Persian
(\fa{۳۲}), but arithmetic values such as \code{15} or loop bounds such as \code{5}
are written with the ordinary digits \code{0}--\code{9}. The Persian digits are part
of the \emph{spelling} of certain keywords; computation uses the standard digits.
\end{note}

\section{Why bilingualism matters}

This feature is more than a novelty. For a learner whose first language is
Persian, writing \code{\fa{اگر}} instead of \code{if} removes a layer of
translation between thought and code --- the program reads the way they think. A
classroom can teach loops and functions without first teaching English. And
because the structure is identical, the day that learner needs to work in English
--- to use an English library, join an international team, read English
documentation --- the transition is trivial: they already know every concept, and
only the spellings change.

\begin{tip}
If you are learning, start in whichever language is more comfortable, and once a
concept is solid, try rewriting a program in the other language. Because only the
keywords differ, this is quick --- and it cements both the idea \emph{and} its
English vocabulary at once, which is exactly the bridge Salam was built to be.
\end{tip}

\begin{deep}[In Depth: how the compiler reads two languages]
Internally, Salam keeps a \emph{language pack}: a table mapping each keyword to its
spellings in each supported language. When the lexer (the first stage of the
compiler, from Chapter~1) reads your source, it consults the pack for the language
you selected and recognises \code{\fa{تابع}} or \code{func} as the very same
internal token. From that point on --- parsing, type-checking, code
generation --- the compiler has no idea which language the source was written in,
because the difference has already been erased. That is why the two versions
produce byte-for-byte identical programs: after the first stage, there is only one
language left, the language of meaning.
\end{deep}

\section{Chapter summary}

\begin{itemize}[leftmargin=1.4em]
  \item Salam's grammar is identical in English and Persian; only the spelling of
        keywords and type names differs.
  \item Every keyword has a Persian twin (\code{func}/\code{\fa{تابع}},
        \code{if}/\code{\fa{اگر}}, \dots); the entry point is \code{main} /
        \code{\fa{اصلی}}. The full table is in Appendix~A.
  \item Compile a Persian file with \code{--lang=fa}; English is the default. The
        rest of the workflow is unchanged.
  \item A file uses one keyword language, but identifiers, strings, and comments
        may be in any language.
  \item Persian digits appear inside certain keyword spellings; computation uses
        ordinary digits.
  \item Bilingualism lets learners code in their own language and cross to English
        effortlessly, because the concepts are identical.
\end{itemize}

\section{Exercises}

\exercise Take any English program from an earlier chapter and rewrite it in
Persian, then run it with \code{--lang=fa}. Confirm the output matches.

\exercise Write a Persian program with a function (\code{\fa{تابع}}) that returns
(\code{\fa{بازگشت}}) the larger of two numbers, using \code{\fa{اگر}}.

\exercise In an English program, use a Persian word (in Latin letters, like
\code{naam}) as a variable name and a Persian sentence as a printed string.

\exercise Write the same ``count to ten'' loop twice --- once in English, once in
Persian --- and place them side by side to see how only the keywords change.

\exercise Explain in a comment why two versions of a program, one English and one
Persian, compile to identical results.
